\documentclass[runningheads]{llncs}

\usepackage{graphicx}
\usepackage{listings}
\usepackage{xcolor}
\usepackage{hyperref}
\usepackage{booktabs}
\usepackage{multirow}
\usepackage{amsmath}
\usepackage{algorithm}
\usepackage{algpseudocode}

% Code listing settings
\lstset{
    basicstyle=\ttfamily\small,
    breaklines=true,
    frame=single,
    numbers=left,
    numberstyle=\tiny,
    keywordstyle=\color{blue},
    commentstyle=\color{gray},
    stringstyle=\color{red},
    showstringspaces=false
}

% Define JSON language for listings
\lstdefinelanguage{json}{
    basicstyle=\ttfamily\small,
    numbers=left,
    numberstyle=\tiny,
    stepnumber=1,
    numbersep=8pt,
    showstringspaces=false,
    breaklines=true,
    frame=single,
    string=[s]{"}{"},
    stringstyle=\color{red},
    comment=[l]{//},
    commentstyle=\color{gray}\ttfamily,
    morestring=[b]',
    morestring=[b]",
    keywordstyle=\color{blue},
    keywords={true,false,null}
}

% Title information
\title{Secure Chat System: Achieving CIANR through Cryptographic Primitives}
\subtitle{Assignment \#2 - Information Security (CS-3002)}

\author{Umer Farooq\inst{1}}
\institute{National University of Computer and Emerging Sciences (FAST--NUCES)\\
\email{i220891@nu.edu.pk}\\
Roll No: 22I-0891, Section: CS-7D\\
Instructor: Urooj Ghani}

\begin{document}

\maketitle

\begin{abstract}
This paper presents the design and implementation of a secure client-server chat system that achieves Confidentiality, Integrity, Authenticity, and Non-Repudiation (CIANR) through the integration of cryptographic primitives. The system implements a self-built Public Key Infrastructure (PKI) with X.509 certificates, Diffie-Hellman key exchange, AES-128 encryption, RSA signatures, and SHA-256 hashing. The protocol ensures secure credential transmission, encrypted message exchange, and provides cryptographic proof of communication through signed session transcripts. Comprehensive testing demonstrates 100\% success rate across 88 automated test cases, validating the correctness of the implementation.

\keywords{Information Security \and Cryptography \and PKI \and CIANR \and Secure Communication}
\end{abstract}

\section{Introduction}

Modern secure communication systems must protect against various threats including passive eavesdropping, active man-in-the-middle attacks, and unauthorized access attempts. This project implements a console-based secure chat system that demonstrates how cryptographic primitives can be combined to achieve comprehensive security properties.

\subsection{Threat Model}

The system is designed to protect against:
\begin{itemize}
    \item \textbf{Passive Eavesdroppers:} Adversaries who can observe network traffic
    \item \textbf{Active MitM Attackers:} Adversaries who can intercept, modify, replay, or inject messages
    \item \textbf{Untrusted Clients:} Attempts at unauthorized login or password guessing
\end{itemize}

\subsection{Security Goals}

The system achieves the following security properties:
\begin{itemize}
    \item \textbf{Confidentiality:} All messages encrypted with AES-128
    \item \textbf{Integrity:} SHA-256 hashing and RSA signatures prevent undetected tampering
    \item \textbf{Authenticity:} X.509 certificate-based mutual authentication
    \item \textbf{Non-Repudiation:} Signed session transcripts provide cryptographic proof
\end{itemize}

\subsection{System Overview}

The secure chat system consists of:
\begin{itemize}
    \item A Python-based client-server architecture
    \item Self-built Root Certificate Authority (CA)
    \item X.509 certificate issuance and validation
    \item MySQL database for secure credential storage
    \item Application-layer cryptographic protocol (no TLS/SSL)
\end{itemize}

\section{Secure Chat Protocol}

The protocol is divided into four major phases that collectively ensure CIANR throughout the communication lifecycle.

\subsection{Control Plane: Negotiation and Authentication}

The control plane establishes trust and identity before any sensitive data exchange. Both parties present X.509 certificates issued by the same trusted CA and perform mutual verification.

\subsubsection{Certificate Exchange}

The client initiates connection by sending a \texttt{HelloMessage} containing:
\begin{itemize}
    \item Client certificate (PEM format)
    \item Nonce for freshness
\end{itemize}

The server responds with a \texttt{ServerHelloMessage} containing:
\begin{itemize}
    \item Server certificate (PEM format)
    \item Nonce for freshness
\end{itemize}

Both parties validate received certificates by checking:
\begin{enumerate}
    \item Signature chain validity (certificate signed by trusted CA)
    \item Expiry date and validity period
    \item Common Name (CN) or hostname match
\end{enumerate}

Invalid certificates result in \texttt{BAD\_CERT} error and connection termination.

\subsubsection{Credential Encryption}

A temporary Diffie-Hellman exchange establishes an ephemeral AES key for credential encryption:

\begin{enumerate}
    \item Client generates DH parameters $(p, g)$ and public value $A = g^a \bmod p$
    \item Server responds with public value $B = g^b \bmod p$
    \item Both compute shared secret: $K_s = A^b \bmod p = B^a \bmod p$
    \item Temporary AES key: $K_{temp} = \text{Trunc}_{16}(\text{SHA256}(\text{big-endian}(K_s)))$
\end{enumerate}

Registration and login credentials are encrypted using $K_{temp}$ before transmission.

\subsubsection{Message Formats}

\textbf{Hello Message:}
\begin{lstlisting}[language=json]
{
  "type": "hello",
  "client_cert": "...PEM...",
  "nonce": "base64_encoded_nonce"
}
\end{lstlisting}

\textbf{Server Hello Message:}
\begin{lstlisting}[language=json]
{
  "type": "server_hello",
  "server_cert": "...PEM...",
  "nonce": "base64_encoded_nonce"
}
\end{lstlisting}

\textbf{Register Message:}
\begin{lstlisting}[language=json]
{
  "type": "register",
  "email": "user@example.com",
  "username": "alice",
  "pwd": "base64_encrypted_password",
  "salt": "base64_salt"
}
\end{lstlisting}

\textbf{Login Message:}
\begin{lstlisting}[language=json]
{
  "type": "login",
  "email": "user@example.com",
  "pwd": "base64_encrypted_password",
  "nonce": "base64_nonce"
}
\end{lstlisting}

\subsection{Key Agreement: Post-Authentication}

After successful authentication, both sides establish a new session key for chat encryption using a fresh Diffie-Hellman exchange.

\subsubsection{Session Key Derivation}

\begin{enumerate}
    \item Client generates session DH parameters $(p, g)$ and public value $A = g^a \bmod p$
    \item Server responds with public value $B = g^b \bmod p$
    \item Both compute shared secret: $K_s = A^b \bmod p = B^a \bmod p$
    \item Session key: $K = \text{Trunc}_{16}(\text{SHA256}(\text{big-endian}(K_s)))$
\end{enumerate}

The 16-byte key $K$ is used for AES-128 encryption throughout the chat session.

\subsubsection{Message Formats}

\textbf{DH Client Message:}
\begin{lstlisting}[language=json]
{
  "type": "dh_client",
  "g": 2,
  "p": 1234567890...,
  "A": 9876543210...
}
\end{lstlisting}

\textbf{DH Server Message:}
\begin{lstlisting}[language=json]
{
  "type": "dh_server",
  "B": 1122334455...
}
\end{lstlisting}

\subsection{Data Plane: Encrypted Message Exchange}

The data plane applies multiple layers of protection: encryption for confidentiality, hashing and digital signatures for integrity and authenticity, and sequence numbers for replay protection.

\subsubsection{Message Transmission Process}

For each message:
\begin{enumerate}
    \item Plaintext is padded using PKCS\#7
    \item Message encrypted: $ct = \text{AES-128-ECB}(plaintext, K)$
    \item Hash computed: $h = \text{SHA256}(seqno \| timestamp \| ct)$
    \item Signature created: $sig = \text{RSA-SIGN}(h, private\_key)$
    \item Message sent as JSON with all fields
\end{enumerate}

\subsubsection{Message Reception Process}

Upon receiving a message:
\begin{enumerate}
    \item Sequence number validated (strictly increasing)
    \item Hash recomputed: $h' = \text{SHA256}(seqno \| timestamp \| ct)$
    \item Signature verified: $\text{RSA-VERIFY}(sig, h', public\_key)$
    \item If valid, decrypt: $plaintext = \text{AES-128-ECB-DECRYPT}(ct, K)$
    \item PKCS\#7 padding removed
\end{enumerate}

\subsubsection{Message Format}

\textbf{Chat Message:}
\begin{lstlisting}[language=json]
{
  "type": "msg",
  "seqno": 1,
  "ts": 1699123456789,
  "ct": "base64_encrypted_ciphertext",
  "sig": "base64_rsa_signature"
}
\end{lstlisting}

Where:
\begin{itemize}
    \item \texttt{seqno}: Sequence number (prevents replay attacks)
    \item \texttt{ts}: Unix timestamp in milliseconds (prevents stale messages)
    \item \texttt{ct}: Base64-encoded AES-128 ciphertext
    \item \texttt{sig}: Base64-encoded RSA signature over $\text{SHA256}(seqno \| ts \| ct)$
\end{itemize}

\subsection{Non-Repudiation: Session Evidence}

The final stage ensures neither participant can deny sending or receiving messages through cryptographic proof.

\subsubsection{Transcript Management}

Each side maintains an append-only transcript file containing:
\begin{lstlisting}
seqno | timestamp | ciphertext | signature | peer_cert_fingerprint
\end{lstlisting}

\subsubsection{Session Receipt Generation}

At session end:
\begin{enumerate}
    \item Compute transcript hash: $H = \text{SHA256}(\text{concatenation of all transcript lines})$
    \item Sign hash: $sig = \text{RSA-SIGN}(H, private\_key)$
    \item Generate SessionReceipt
\end{enumerate}

\subsubsection{Session Receipt Format}

\begin{lstlisting}[language=json]
{
  "type": "receipt",
  "peer": "client",
  "first_seq": 1,
  "last_seq": 10,
  "transcript_sha256": "hex_hash_of_transcript",
  "sig": "base64_rsa_signature"
}
\end{lstlisting}

\subsubsection{Offline Verification}

Any third party can verify:
\begin{enumerate}
    \item Recompute transcript hash from transcript file
    \item Verify RSA signature using participant's certificate
    \item Confirm authenticity and completeness
\end{enumerate}

\section{System Requirements and Implementation}

\subsection{PKI Setup and Certificate Validation}

\subsubsection{Root CA Generation}

The system implements \texttt{scripts/gen\_ca.py} to create a self-signed Root CA:
\begin{itemize}
    \item Generates RSA-2048 private key
    \item Creates self-signed X.509 certificate
    \item Stores CA key and certificate in \texttt{certs/} directory
\end{itemize}

\subsubsection{Certificate Issuance}

The system implements \texttt{scripts/gen\_cert.py} to issue certificates:
\begin{itemize}
    \item Generates RSA-2048 keypair for entity (server/client)
    \item Creates X.509 certificate signed by Root CA
    \item Includes Common Name (CN) and Subject Alternative Name (SAN)
    \item Stores certificate and private key in PEM format
\end{itemize}

\subsubsection{Certificate Validation}

The system validates certificates by checking:
\begin{enumerate}
    \item \textbf{Chain Validity:} Certificate signed by trusted CA
    \item \textbf{Expiry:} Certificate not expired
    \item \textbf{Hostname:} CN or SAN matches expected value
\end{enumerate}

Invalid certificates (self-signed, expired, wrong hostname) are rejected with \texttt{BAD\_CERT} error.

\subsubsection{Certificate Inspection}

Certificate details can be inspected using:
\begin{lstlisting}[language=bash]
openssl x509 -in certs/server-cert.pem -text -noout
\end{lstlisting}

This reveals:
\begin{itemize}
    \item Certificate issuer and subject
    \item Validity period
    \item Public key information
    \item Extensions (CN, SAN)
    \item Signature algorithm
\end{itemize}

\subsection{Registration and Login Security}

\subsubsection{Password Storage}

User credentials are stored in MySQL database with:
\begin{itemize}
    \item Per-user random salt (16 bytes minimum)
    \item Salted password hash: $pwd\_hash = \text{hex}(\text{SHA256}(salt \| password))$
    \item No plaintext passwords in database or logs
\end{itemize}

\subsubsection{Database Schema}

\begin{lstlisting}[language=sql]
CREATE TABLE users (
    email VARCHAR(255) PRIMARY KEY,
    username VARCHAR(255) UNIQUE NOT NULL,
    salt VARBINARY(16) NOT NULL,
    pwd_hash CHAR(64) NOT NULL,
    created_at TIMESTAMP DEFAULT CURRENT_TIMESTAMP,
    last_login TIMESTAMP NULL DEFAULT NULL
);
\end{lstlisting}

\subsubsection{Authentication Flow}

\begin{enumerate}
    \item Client connects and exchanges certificates (mutual verification)
    \item Temporary DH exchange establishes $K_{temp}$
    \item Credentials encrypted with $K_{temp}$ and sent to server
    \item Server decrypts and verifies:
    \begin{itemize}
        \item Client certificate is valid and trusted
        \item For registration: email/username not already registered
        \item For login: salted hash matches stored hash
    \end{itemize}
    \item Authentication succeeds only if both conditions met
\end{enumerate}

\subsection{Session Key Establishment}

After successful authentication, a new Diffie-Hellman exchange establishes the chat session key:

\begin{algorithm}
\caption{Session Key Establishment}
\begin{algorithmic}[1]
\State Client generates DH parameters $(p, g)$
\State Client chooses private exponent $a$, computes $A = g^a \bmod p$
\State Client sends $(p, g, A)$ to server
\State Server chooses private exponent $b$, computes $B = g^b \bmod p$
\State Server sends $B$ to client
\State Client computes: $K_s = B^a \bmod p$
\State Server computes: $K_s = A^b \bmod p$
\State Both derive: $K = \text{Trunc}_{16}(\text{SHA256}(\text{big-endian}(K_s)))$
\end{algorithmic}
\end{algorithm}

\subsection{Encrypted Chat and Message Integrity}

\subsubsection{Encryption}

Messages are encrypted using AES-128 in ECB mode with PKCS\#7 padding:
\begin{itemize}
    \item Plaintext padded to 16-byte boundary
    \item Each 16-byte block encrypted independently
    \item Ciphertext base64-encoded for transmission
\end{itemize}

\subsubsection{Integrity and Authenticity}

Each message includes:
\begin{itemize}
    \item SHA-256 hash: $h = \text{SHA256}(seqno \| timestamp \| ciphertext)$
    \item RSA signature: $sig = \text{RSA-SIGN}(h, sender\_private\_key)$
    \item Signature uses PKCS\#1 v1.5 with SHA-256
\end{itemize}

\subsubsection{Replay Protection}

Sequence numbers ensure:
\begin{itemize}
    \item Messages processed in order
    \item Replayed messages rejected (same seqno detected)
    \item Out-of-order messages rejected
\end{itemize}

\subsection{Non-Repudiation and Session Closure}

\subsubsection{Transcript Maintenance}

Each participant maintains an append-only transcript:
\begin{lstlisting}
1|1699123456789|base64_ct|base64_sig|peer_fingerprint
2|1699123456790|base64_ct|base64_sig|peer_fingerprint
...
\end{lstlisting}

\subsubsection{Receipt Generation}

At session end:
\begin{enumerate}
    \item Compute: $H = \text{SHA256}(\text{all transcript lines})$
    \item Sign: $sig = \text{RSA-SIGN}(H, private\_key)$
    \item Generate SessionReceipt JSON file
\end{enumerate}

\subsubsection{Verification}

Offline verification confirms:
\begin{itemize}
    \item Each message signature is valid
    \item Transcript hash matches receipt
    \item Receipt signature is valid
    \item Any modification breaks verification
\end{itemize}

\section{Testing and Evidence}

\subsection{Automated Test Suite}

The system includes a comprehensive automated test suite covering all cryptographic primitives and protocol phases.

\subsubsection{Test Coverage}

The test suite includes 88 test cases across 12 test modules:

\begin{table}[h]
\centering
\begin{tabular}{@{}llr@{}}
\toprule
\textbf{Test Module} & \textbf{Description} & \textbf{Tests} \\
\midrule
\texttt{test\_protocol.py} & Protocol message models & 12 \\
\texttt{test\_cert\_gen.py} & Certificate generation & 6 \\
\texttt{test\_db.py} & Database operations & 9 \\
\texttt{test\_transcript.py} & Transcript management & 9 \\
\texttt{test\_server\_client.py} & Server/Client integration & 3 \\
\texttt{test\_utils.py} & Utility functions & 7 \\
\texttt{test\_aes.py} & AES encryption & 8 \\
\texttt{test\_sign.py} & RSA signatures & 8 \\
\texttt{test\_dh.py} & Diffie-Hellman key exchange & 7 \\
\texttt{test\_pki.py} & PKI certificate validation & 9 \\
\texttt{test\_integration.py} & Crypto integration & 4 \\
\texttt{test\_config.py} & Configuration system & 6 \\
\bottomrule
\end{tabular}
\caption{Test Module Coverage}
\label{tab:test-coverage}
\end{table}

\subsubsection{Test Results}

The automated test suite was executed on November 16, 2025, with the following results:

\begin{table}[h]
\centering
\begin{tabular}{@{}lr@{}}
\toprule
\textbf{Metric} & \textbf{Value} \\
\midrule
Total Tests Run & 88 \\
Passed & 88 \\
Failed & 0 \\
Errors & 0 \\
Skipped & 9 \\
Success Rate & 100.0\% \\
Overall Status & PASS \\
\bottomrule
\end{tabular}
\caption{Automated Test Results}
\label{tab:test-results}
\end{table}

\paragraph{Skipped Tests}

Nine database tests were skipped due to database unavailability during test execution. These tests validate:
\begin{itemize}
    \item Database connection
    \item User registration
    \item User authentication
    \item Password hashing
    \item Duplicate prevention
\end{itemize}

The database functionality is implemented and tested separately on systems with MySQL access.

\subsection{Manual Security Tests}

\subsubsection{Invalid Certificate Tests}

The system rejects invalid certificates with \texttt{BAD\_CERT} errors:
\begin{itemize}
    \item \textbf{Self-signed certificates:} Not signed by trusted CA
    \item \textbf{Expired certificates:} Past validity period
    \item \textbf{Wrong hostname:} CN/SAN mismatch
\end{itemize}

\subsubsection{Message Tampering Detection}

Tampering with ciphertext results in \texttt{SIG\_FAIL} errors:
\begin{itemize}
    \item Modified ciphertext causes hash mismatch
    \item Signature verification fails
    \item Message rejected by recipient
\end{itemize}

\subsubsection{Replay Attack Prevention}

Replayed messages are detected and rejected:
\begin{itemize}
    \item Duplicate sequence numbers trigger \texttt{REPLAY} error
    \item Out-of-order messages rejected
    \item Sequence number tracking maintained correctly
\end{itemize}

\subsection{Non-Repudiation Verification}

\subsubsection{Transcript Verification}

Each message in the transcript can be verified:
\begin{enumerate}
    \item Recompute: $h = \text{SHA256}(seqno \| timestamp \| ciphertext)$
    \item Verify: $\text{RSA-VERIFY}(signature, h, sender\_public\_key)$
    \item Confirm message authenticity
\end{enumerate}

\subsubsection{Receipt Verification}

The SessionReceipt can be verified:
\begin{enumerate}
    \item Recompute transcript hash from transcript file
    \item Verify: $\text{RSA-VERIFY}(receipt\_sig, transcript\_hash, sender\_public\_key)$
    \item Confirm session authenticity and completeness
\end{enumerate}

\subsubsection{Tampering Detection}

Any modification to the transcript:
\begin{itemize}
    \item Changes the transcript hash
    \item Invalidates the receipt signature
    \item Breaks offline verification
\end{itemize}

\section{Implementation Details}

\subsection{Technology Stack}

\begin{itemize}
    \item \textbf{Language:} Python 3.8+
    \item \textbf{Cryptography:} \texttt{cryptography} library (Hazmat primitives)
    \item \textbf{Database:} MySQL 8.0+ / MariaDB 10.3+
    \item \textbf{Protocol:} Plain TCP sockets (no TLS/SSL)
    \item \textbf{Serialization:} JSON with Pydantic models
\end{itemize}

\subsection{Key Implementation Files}

\begin{itemize}
    \item \texttt{scripts/gen\_ca.py}: Root CA generation
    \item \texttt{scripts/gen\_cert.py}: Certificate issuance
    \item \texttt{app/crypto/pki.py}: Certificate validation
    \item \texttt{app/crypto/dh.py}: Diffie-Hellman key exchange
    \item \texttt{app/crypto/aes.py}: AES-128 encryption
    \item \texttt{app/crypto/sign.py}: RSA signatures
    \item \texttt{app/storage/db.py}: Database operations
    \item \texttt{app/storage/transcript.py}: Transcript management
    \item \texttt{app/server.py}: Server implementation
    \item \texttt{app/client.py}: Client implementation
\end{itemize}

\subsection{Security Considerations}

\subsubsection{Key Management}

\begin{itemize}
    \item Private keys stored locally, never transmitted
    \item CA private key never committed to version control
    \item Session keys derived from ephemeral DH exchange
    \item Each session uses unique key
\end{itemize}

\subsubsection{Password Security}

\begin{itemize}
    \item Per-user random salt (16 bytes)
    \item SHA-256 hashing: $hash = \text{SHA256}(salt \| password)$
    \item Constant-time comparison to prevent timing attacks
    \item No plaintext passwords in database or logs
\end{itemize}

\subsubsection{Protocol Security}

\begin{itemize}
    \item Certificate validation before credential exchange
    \item Encrypted credential transmission
    \item Message integrity through signatures
    \item Replay protection via sequence numbers
    \item Timestamp validation for freshness
\end{itemize}

\section{Certificate Inspection Results}

\subsection{Root CA Certificate}

The Root CA certificate includes:
\begin{itemize}
    \item \textbf{Issuer:} FAST-NU Root CA (self-signed)
    \item \textbf{Subject:} FAST-NU Root CA
    \item \textbf{Validity:} 10 years from generation
    \item \textbf{Key Algorithm:} RSA 2048-bit
    \item \textbf{Signature Algorithm:} SHA256 with RSA
    \item \textbf{Extensions:} Basic Constraints (CA:TRUE)
\end{itemize}

\subsection{Server Certificate}

The server certificate includes:
\begin{itemize}
    \item \textbf{Issuer:} FAST-NU Root CA
    \item \textbf{Subject CN:} server.local
    \item \textbf{Subject Alternative Name:} DNS:server.local
    \item \textbf{Validity:} 1 year from issuance
    \item \textbf{Key Algorithm:} RSA 2048-bit
    \item \textbf{Signature Algorithm:} SHA256 with RSA
\end{itemize}

\subsection{Client Certificate}

The client certificate includes:
\begin{itemize}
    \item \textbf{Issuer:} FAST-NU Root CA
    \item \textbf{Subject CN:} client.local
    \item \textbf{Subject Alternative Name:} DNS:client.local
    \item \textbf{Validity:} 1 year from issuance
    \item \textbf{Key Algorithm:} RSA 2048-bit
    \item \textbf{Signature Algorithm:} SHA256 with RSA
\end{itemize}

\section{Security Analysis}

\subsection{Confidentiality}

\begin{itemize}
    \item \textbf{AES-128 Encryption:} All messages encrypted with session key
    \item \textbf{Key Derivation:} Session key from DH exchange (forward secrecy)
    \item \textbf{Credential Protection:} Passwords encrypted before transmission
    \item \textbf{No Plaintext:} No sensitive data transmitted in clear text
\end{itemize}

\subsection{Integrity}

\begin{itemize}
    \item \textbf{SHA-256 Hashing:} Message digest prevents tampering
    \item \textbf{RSA Signatures:} Cryptographic signatures ensure authenticity
    \item \textbf{Tampering Detection:} Any modification invalidates signature
    \item \textbf{Replay Protection:} Sequence numbers prevent message replay
\end{itemize}

\subsection{Authenticity}

\begin{itemize}
    \item \textbf{Mutual Authentication:} Both client and server verify certificates
    \item \textbf{Certificate Validation:} Chain, expiry, and hostname checks
    \item \textbf{Signature Verification:} Each message signed and verified
    \item \textbf{Identity Binding:} Certificates bind public keys to identities
\end{itemize}

\subsection{Non-Repudiation}

\begin{itemize}
    \item \textbf{Transcript Logging:} Append-only transcript of all messages
    \item \textbf{Signed Receipt:} Transcript hash signed by participant
    \item \textbf{Offline Verification:} Third parties can verify authenticity
    \item \textbf{Tamper Evidence:} Any modification breaks verification
\end{itemize}

\section{Conclusion}

This project successfully implements a secure chat system that achieves CIANR through the integration of cryptographic primitives. The system demonstrates:

\begin{itemize}
    \item Proper PKI setup with self-built CA and certificate issuance
    \item Secure credential handling with salted hashing and encrypted transmission
    \item Authenticated key agreement using Diffie-Hellman
    \item Encrypted message exchange with integrity protection
    \item Non-repudiation through signed session transcripts
\end{itemize}

Comprehensive testing validates the correctness of the implementation, with 88 automated tests achieving 100\% pass rate. The system provides a practical demonstration of how cryptographic primitives combine to achieve comprehensive security in real-world applications.

\section*{Acknowledgments}

The implementation follows the assignment specification provided for CS-3002 Information Security course at FAST--NUCES. Cryptographic primitives are implemented using the Python \texttt{cryptography} library.


\end{document}

